% Recommended Sidechain Compressor settings, derived from a 150-point
% Latin Hypercube parameter sweep (analysis/run_lhs_sweep.py). See
% analysis/lhs_sweep_findings.txt for full methodology and the 4 findings
% this table summarizes.
%
% Requires packages already in main.tex's preamble: booktabs, siunitx.
% Also requires colortbl for \cellcolor - see parameter_sweep_table.tex's
% header comment (add \usepackage{colortbl} or change to
% \usepackage[table]{xcolor}).
%
% KEY FINDING this table is built on: mean margin gain, STOI gain, and
% psychoacoustic partial-loudness reduction are almost perfectly
% correlated with each other across the swept space (r=0.994 for both
% STOI and loudness reduction against margin gain) - there is no
% parameter combination that decouples them. The real tradeoff is
% threshold vs. how much of the clip the sidechain actually engages on
% (measured via the dry Dialogue stem's own envelope crossing each
% threshold), which directly opposes wanting the mix untouched when
% Dialogue isn't present. These three rows are real tested points from
% the sweep spanning that tradeoff, not a single "optimal" answer -
% ratio/knee/attack/release matter far less than threshold (all
% |r|<0.19 individually) so are not the load-bearing part of any row.

\definecolor{paramgrey}{gray}{0.85}

\begin{table}[h]
\centering
\footnotesize
\caption{Recommended Sidechain Compressor operating points, Advanced
(Summed-bus) ducking, spanning the threshold/engagement tradeoff
identified in the 150-point parameter sweep. ``Engaged'' is the fraction
of the clip where the dry Dialogue stem's own smoothed envelope exceeds
the threshold (i.e. where ducking of the other channels can occur at
all). Gains are relative to the unprocessed baseline.}
\label{tab:recommended-settings}
\begin{tabular}{@{}l S[table-format=-2.1] S[table-format=1.1] S[table-format=2.1] S[table-format=2.1] S[table-format=3.0] S[table-format=2.1] S[table-format=+1.2] S[table-format=+1.3] S[table-format=+2.1]@{}}
\toprule
\cellcolor{paramgrey}\textbf{Tier} & \cellcolor{paramgrey}{\textbf{Thr.} (\si{\decibel})} & \cellcolor{paramgrey}{\textbf{Ratio}} & \cellcolor{paramgrey}{\textbf{Knee} (\si{\decibel})} & \cellcolor{paramgrey}{\textbf{Atk.} (\si{\milli\second})} & \cellcolor{paramgrey}{\textbf{Rel.} (\si{\milli\second})} & \cellcolor{paramgrey}{\textbf{Engaged} (\si{\percent})} & \textbf{Margin Gain} (\si{\decibel}) & \textbf{STOI Gain} & \textbf{Loud.\ Red.\ (pp)} \\
\midrule
Aggressive   & -58.6 & 7.4 & 5.9  & 3.2  & 215 & 68.8 & \cellcolor{green!70}+6.96 & \cellcolor{yellow!70}+0.104 & \cellcolor{orange!70}+11.4 \\
Balanced     & -40.3 & 5.7 & 22.9 & 16.9 & 347 & 29.9 & \cellcolor{green!60}+6.21 & \cellcolor{yellow!67}+0.100 & \cellcolor{orange!60}+9.8  \\
Conservative & -28.8 & 7.7 & 9.4  & 8.5  & 384 &  1.9 & \cellcolor{green!23}+2.40 & \cellcolor{yellow!27}+0.042 & \cellcolor{orange!25}+3.0  \\
\bottomrule
\end{tabular}
\par\vspace{4pt}
\raggedright\footnotesize Threshold is the dominant parameter
($r=-0.929$ with margin gain across the sweep); ratio, knee, attack, and
release each correlate weakly on their own ($|r|<0.19$), so the exact
decimal values in those four columns are not individually load-bearing --
what matters is the threshold tier. The sweep's search range
(\SI{-60}{\decibel} to \SI{-18}{\decibel}) was still improving at its
lower edge, so an even more aggressive tier below \SI{-60}{\decibel}
was not ruled out, only untested.
\end{table}
